\ncdmodultest{2}{Ferie, Arbejdsliv og Fritid}{全真模考卷二：夏日住宿、职场面试与农园小屋}{丹麦夏日与周末度假住宿 / 欧登塞面试路上的意外 / 市民农园小木屋文化 / 三人谈工作时间与职场生活 / 新兼职回信与物业报修 / 语言学校与跑步俱乐部口语及大学生信息差}{形容词比较级与最高级（-ere/-st, flere/mere）·过去完成时 havde/var + 分词 · 时间介词 i / for ... siden / om / indtil · 礼貌半正式报修沟通}

% =====================================================================
% Del 1: Læsning
% =====================================================================
\ncdsubsec{Del 1: Læsning（阅读理解 · 40 minutter · 24 point）}

\begin{tipbox}{考场规则与配速建议（Officielle regler）}
\small
\textit{Opgave 1, 2, 3 og 4 hæftes og udleveres samlet. Der afsættes \textbf{40 minutter} til de fire opgaver i alt. Kursisten må ikke bruge hjælpemidler.}（阅读四题合订发放，\textbf{总计时 40 分钟}，禁用任何词典或手机。及格线：\textbf{16--24 point = Bestået}；0--15 point = Ikke bestået。建议时间分配：Opgave 1 (8 min) + Opgave 2 (8 min) + Opgave 3 (8 min) + Opgave 4 (12 min) + 复查 (4 min)。）
\end{tipbox}

% ---------------------------------------------------------------------
\subsection*{Opgave 1 (NY · 6 point) --- \textit{Sommerferie og overnatning i Danmark}}

\textbf{Instruktion:}
\begin{itemize}
  \item Læs de seks små tekster om personerne, der gerne vil på ferie eller weekendtur i Danmark (1--6).
  \item Find den annonce (B--I), der passer til hver tekst.
  \item Skriv bogstavet på annoncen, der passer til teksten. Der er to annoncer, du ikke skal bruge. Se eksemplet (0).
\end{itemize}

\begin{tcolorbox}[enhanced, breakable, colback=NCDBgGray, colframe=NCDGray, coltitle=white, fonttitle=\bfseries, title={Personerne (0--6)}, arc=2pt, boxrule=0.8pt, borderline west={3pt}{0pt}{NCDNavy}, left=10pt, right=10pt, top=7pt, bottom=7pt]
\small
\textbf{0.} Mette og Henrik har to små børn og en stor hund. De vil gerne leje et sommerhus tæt på stranden i Vestjylland i en uge. Det er vigtigt for dem, at de selv kan lave mad, og at hunden må være med udenfor i en have. \quad \textcolor{NCDRed}{\textbf{[Eksempel: Svar A]}}

\vspace{0.4em}
\textbf{1.} Nadia og hendes søn på fem år elsker dyr og vil gerne prøve en anderledes sommerferie i fem dage. Drengen drømmer om at se geder eller høns tæt på og hjælpe med at samle æg om morgenen. De har deres egen bil, så stedet behøver ikke at ligge i en by. \quad \textbf{Svar:} \rule{1.2cm}{0.4pt}

\vspace{0.4em}
\textbf{2.} Mogens og Pia på 58 år vil holde fire dages afslappende ferie på en dansk ø i sommer. De gider ikke lave mad selv på ferien, så de vil gerne bo et sted med både morgenmad og aftensmad. De elsker at cykle stille ture rundt i et smukt landskab. \quad \textbf{Svar:} \rule{1.2cm}{0.4pt}

\vspace{0.4em}
\textbf{3.} Kasper og hans to studiekammerater på 22 år vil på en billig weekendtur til Nordjylland for at surfe og møde andre unge. De har ikke bil, så de rejser med tog og bus. Derfor skal der være let adgang til offentlig transport og et fælleskøkken, hvor de kan lave mad. \quad \textbf{Svar:} \rule{1.2cm}{0.4pt}

\vspace{0.4em}
\textbf{4.} Daniel er journalist og vil tage fem dage væk alene for at skrive på sin nye bog i rolige omgivelser ved havet. Han har brug for et lille, fredeligt sted med et godt skrivebord og hurtigt internet, så han kan arbejde om formiddagen og gå ture ved vandet om eftermiddagen. \quad \textbf{Svar:} \rule{1.2cm}{0.4pt}

\vspace{0.4em}
\textbf{5.} Birgit på 68 år rejser med sin mand, der bruger kørestol efter en operation. De vil på en forlænget weekend til en større by, hvor de let kan gå på museum og spise på restaurant uden at bruge bil. Deres bolig skal ligge tæt på stationen og have elevator. \quad \textbf{Svar:} \rule{1.2cm}{0.4pt}

\vspace{0.4em}
\textbf{6.} Line og Thomas trænger til en rolig weekend kun for to, helt uden skærme, arbejde og internet. De elsker at sove ude i naturen i smukke omgivelser, men de vil ikke ligge på et hårdt liggeunderlag. De vil have en rigtig seng og morgenmad bragt til døren. \quad \textbf{Svar:} \rule{1.2cm}{0.4pt}
\end{tcolorbox}

\vspace{0.3em}
\textbf{Annoncer om overnatning (A--I):}

{\small
\renewcommand{\arraystretch}{1.15}
\begin{longtable}{@{} >{\raggedright\arraybackslash}p{8.2cm} @{\hspace{0.6cm}\extracolsep{\fill}} >{\raggedright\arraybackslash}p{8.2cm} @{}}
\toprule
\multicolumn{2}{@{} >{\raggedright\arraybackslash}p{17.0cm} @{}}{%
  \textbf{\color{NCDRed}A. Vestkysten Sommerhusudlejning --- Familievenligt sommerhus i Blåvand \textit{[Eksempel]}}\newline
  Lej et dejligt sommerhus kun 400 meter fra stranden ved Blåvand i Vestjylland! Huset har tre soveværelser, et fuldt udstyret køkken og en stor, indhegnet have, hvor børn og hunde kan lege frit. Her bor du fredeligt tæt på Vesterhavet og kan lave aftensmad på terrassen. Pris: 5.400 kr. for en hel uge. Se flere billeder på \textit{vestkysthus.dk}.
} \\
\midrule
\endfirsthead
\toprule
\multicolumn{2}{@{} >{\raggedright\arraybackslash}p{17.0cm} @{}}{%
  \textbf{\color{NCDGray}Annoncer (fortsat)}
} \\
\midrule
\endhead
\bottomrule
\endfoot
\textbf{\color{NCDNavy}B. Klitmøller Vandrerhjem}\newline
Billig overnatning i Thy i Nordjylland, kun ti minutters gang fra bussen og et af Danmarks bedste surfsteder! Vi tilbyder enkle værelser til unge og aktive gæster fra 240 kr. pr. nat. Her er et stort fælleskøkken og en hyggelig stue, hvor man hurtigt møder andre rejsende. Husdyr er ikke tilladt. Book din seng på \textit{klitmøllerhostel.dk}.
&
\textbf{\color{NCDNavy}F. Solsikken Familiecamping}\newline
Tag hele familien med på en livlig campingferie ved Marielyst Strand på Falster! Medbring dit eget telt eller campingvogn. Pladsen har stor udendørs swimmingpool, legeplads og børnedisko hver aften i juli. Der er hurtigt wi-fi overalt og en lille kiosk, men intet fælleskøkken til gæster. Hunde er velkomne på hele pladsen. Bestil plads på \textit{solsikken-camp.dk}.
\\
\midrule
\textbf{\color{NCDNavy}C. Lilleø Naturshelters \& Glamping}\newline
Trænger I til ro væk fra telefoner og e-mails? Vores luksustelte ligger midt i en stille skov kun for voksne par (ingen børn eller kæledyr). I sover i rigtige senge med varme dyner, og hver morgen leverer vi en kurv med frisk morgenmad til teltet. Bemærk: Der er hverken wi-fi eller tv på pladsen. Se mere på \textit{lilleøglamping.dk}.
&
\textbf{\color{NCDNavy}G. Ærø Badehotel \& Cykelferie}\newline
Nyd et roligt ophold på vores klassiske badehotel på den smukke ø Ærø! Prisen inkluderer både morgenbuffet og en lækker toretters aftensmad i vores restaurant hver dag, plus gratis lån af elcykler under hele ferien. Perfekt til voksne, der vil koble af. Værelserne ligger på 1. sal i en gammel bygning uden elevator. Book 2--7 nætter på \textit{ærøbadehotel.dk}.
\\
\midrule
\textbf{\color{NCDNavy}D. Sønderby Bondegårdsferie}\newline
Oplev livet på landet i en hyggelig ferielejlighed på vores økologiske gård på Fyn! Hver morgen kl. 8.30 må børnene hjælpe med at fodre geder, kaniner og høns samt samle friske æg til morgenmaden. Lejligheden har eget tekøkken og stor legeplads. Der kører ingen bus hertil, så I skal have egen bil. Book ophold (3--7 dage) på \textit{sønderbygård.dk}.
&
\textbf{\color{NCDNavy}H. Skovhytten ved Silkeborg Søerne}\newline
Stor, rustik træhytte til den store familie eller vennegruppe på 6--10 personer! Hytten ligger smukt ved skoven nær Silkeborg og udlejes kun hele uger (minimum syv nætter fra lørdag til lørdag). Der er fire soveværelser, brændeovn og adgang til kanoer, men ingen forplejning og stejle trapper op til hemsen. Bil anbefales. Se ledige uger på \textit{silkeborghytte.dk}.
\\
\midrule
\textbf{\color{NCDNavy}E. Byhavnen Lejlighedshotel}\newline
Komfortable ferielejligheder midt i Aarhus, kun fem minutters gang fra hovedbanegården! Alle etager har bred elevator, og flere lejligheder er indrettet til kørestolsbrugere med ekstra plads på badeværelset. Museer, caféer og restauranter ligger lige uden for døren. Lejlighederne har eget køkken, men du kan også tilkøbe morgenmad. Kæledyr er desværre ikke tilladt. Læs mere på \textit{byhavnen-aarhus.dk}.
&
\textbf{\color{NCDNavy}I. Bornholm Kysthytter}\newline
Små, rolige træhytter til én eller to gæster direkte ved Bornholms klipper og hav! Hver hytte har et lille tekøkken, et godt arbejdsbord og lynhurtigt fibernet, så du kan arbejde uforstyrret eller slappe af ved vandet. Der er god busforbindelse fra Rønne havn og åbent hele året (3--14 nætter). Ingen restaurant og ingen husdyr. Se priser på \textit{bornholmkyst.dk}.
\\
\end{longtable}
}

% ---------------------------------------------------------------------
\subsection*{Opgave 2 (5 point) --- \textit{Anders skal til jobsamtale i Odense}}

\textbf{Instruktion:} Læs teksten \textit{Anders skal til jobsamtale i Odense}. I hvert afsnit er der én sætning, der ikke passer i afsnittet. Find den sætning, der ikke passer. Streg sætningen ud. Se eksemplet i afsnit (0).

\begin{emailbox}{Anders skal til jobsamtale i Odense}
\textbf{(0)} Tirsdag morgen stod Anders meget tidligt op, fordi han skulle til jobsamtale hos et stort firma i Odense. Han havde allerede strøget sin pæneste skjorte aftenen før og pakket sit cv i en sort mappe. \textcolor{NCDRed}{\textbf{[Eksempel --- streges ud:]}}\ \textit{Anders gik direkte i seng igen og sov indtil klokken tre om eftermiddagen.} Derefter spiste han en hurtig morgenmad og cyklede hen til Aarhus Hovedbanegård i god tid.

\vspace{0.35em}
\textbf{(1)} Da han kom hen til perronen, stod der allerede mange travle pendlere i morgenkulden. Anders havde købt sin billet for tre dage siden, så han fandt hurtigt en plads i stillezonen. Toget til Odense kørte præcis klokken 07.15. Det var midt om natten, og stationen var helt tom. Han satte sig ved vinduet og læste sine noter til samtalen en sidste gang.

\vspace{0.35em}
\textbf{(2)} Halvvejs på turen opdagede Anders pludselig, at han havde glemt sin bærbare computer på køkkenbordet derhjemme. Han havde ellers arbejdet på en vigtig præsentation til samtalen indtil sent aftenen før. Anders blev meget nervøs, fordi alle filerne kun lå på den glemte computer. Derfor åbnede han straks sin egen computer på bordet og rettede de sidste billeder. Heldigvis fik han hurtigt fat i sin veninde i Aarhus, som sendte præsentationen til hans mail.

\vspace{0.35em}
\textbf{(3)} Da toget ankom til Odense Banegård, regnede det pludselig kraftigt. Anders havde ikke taget en paraply med, så han ville hellere tage bussen end at gå tyve minutter i regnen. Ved stoppestedet hjalp en venlig dame ham med at finde den hurtigste linje mod erhvervsparken. Solen skinnede fra en helt blå himmel, og alle gik rundt i T-shirt. Anders steg ind i bussen og nåede frem til firmaet ti minutter før mødet.

\vspace{0.35em}
\textbf{(4)} I receptionen blev Anders taget godt imod af chefen, Sofie, som tilbød ham en kop kaffe. Hun havde allerede læst hans ansøgning og syntes, at hans tidligere erfaring var meget relevant. Efter fem minutter gik de ind i et lyst mødelokale, hvor en kollega også sad klar. Sofie gik straks sin vej og lukkede døren, fordi hun aldrig holdt jobsamtaler om tirsdagen. Anders fik hurtigt åbnet sin præsentation på skærmen, og samtalen gik bedre, end han havde turdet håbe.

\vspace{0.35em}
\textbf{(5)} Da mødet var slut, takkede Sofie for en god samtale og lovede at ringe inden fredag middag. På vej tilbage til stationen var Anders ikke længere nervøs, men træt og meget lettet. Han sendte straks en glad besked til sin veninde for at takke for hjælpen tidligere på dagen. Han var så vred over den dårlige samtale, at han aldrig ville besøge Odense igen. Da han sad i toget mod Aarhus, købte han en croissant og nød den stille tur hjem.
\end{emailbox}

% ---------------------------------------------------------------------
\subsection*{Opgave 3 (6 point) --- \textit{Kolonihaver --- danskernes grønne oaser}}

\textbf{Instruktion:} Læs teksten \textit{Kolonihaver --- danskernes grønne oaser}. Skriv de ord, der mangler. Du skal kun bruge hvert ord én gang. Der er fire ord, du ikke skal bruge. Se eksemplet i første linje.

\begin{emailbox}{Kolonihaver --- danskernes grønne oaser}
I udkanten af mange danske byer ligger der grønne haver med små træhuse i \underline{\textbf{alle}} farver. Det er Danmarks berømte kolonihaver, som har eksisteret i over 100 år. I begyndelsen var haverne tænkt til arbejdere i byerne, så de kunne dyrke deres (1)\ \rule{2.1cm}{0.4pt} grøntsager og få frisk luft.

\vspace{0.4em}
I dag er kolonihaver blevet mere populære (2)\ \rule{2.1cm}{0.4pt} nogensinde før, især blandt unge familier i små lejligheder uden altan. Om sommeren tilbringer ejerne næsten hver weekend i haven, (3)\ \rule{2.1cm}{0.4pt} de griller, læser og snakker med naboerne. Mange foreninger har strenge regler. Man må kun (4)\ \rule{2.1cm}{0.4pt} i sit hus (5)\ \rule{2.1cm}{0.4pt} slutningen af oktober, fordi husene ikke er bygget til en kold dansk vinter. Resten af året står de tomme, men (6)\ \rule{2.1cm}{0.4pt} foråret kommer tilbage, fyldes haverne igen med blomster, børn og kaffeduft.
\end{emailbox}

\begin{tcolorbox}[enhanced, colback=NCDBgBlue, colframe=NCDNavy, arc=2pt, boxrule=0.8pt, left=10pt, right=10pt, top=7pt, bottom=7pt]
\centering
\textbf{Ord til opgaven (11 ord: 1 eksempel + 6 svar + 4 ekstra ord):}\\[4pt]
\large
\textit{havde \quad indtil \quad aldrig \quad egne \quad penge \quad når \quad alle \quad sælge \quad bo \quad end \quad hvor}
\end{tcolorbox}

% ---------------------------------------------------------------------
\clearpage
\subsection*{Opgave 4 (NY · 7 point) --- \textit{Tre personer fortæller om deres arbejdsliv og arbejdstider}}

\textbf{Instruktion:}
\begin{itemize}
  \item Læs de tre tekster, og læs spørgsmålene (1--7).
  \item Find den person, der passer til hvert af de syv spørgsmål.
  \item Sæt kryds ud for den person (A, B eller C), der passer til spørgsmålet. Se eksemplet (0).
\end{itemize}

\begin{center}
\renewcommand{\arraystretch}{1.20}
\begin{tabularx}{\linewidth}{@{} >{\raggedright\arraybackslash}X >{\centering\arraybackslash}p{1.9cm} >{\centering\arraybackslash}p{1.9cm} >{\centering\arraybackslash}p{1.9cm} @{}}
\toprule
 & \textbf{A. Freja} & \textbf{B. Karim} & \textbf{C. Lotte} \\
\midrule
\textbf{0.} Hvem arbejder på et hospital? \textit{[Eksempel]} & A\ \ \textbf{X} & B\ \ \underline{\hspace{0.5cm}} & C\ \ \underline{\hspace{0.5cm}} \\
\midrule
\textbf{1.} Hvem kan arbejde hjemme nogle dage om ugen? & A\ \ \underline{\hspace{0.5cm}} & B\ \ \underline{\hspace{0.5cm}} & C\ \ \underline{\hspace{0.5cm}} \\
\textbf{2.} Hvem arbejder på deltid? & A\ \ \underline{\hspace{0.5cm}} & B\ \ \underline{\hspace{0.5cm}} & C\ \ \underline{\hspace{0.5cm}} \\
\textbf{3.} Hvem har arbejdet det samme sted i længst tid? & A\ \ \underline{\hspace{0.5cm}} & B\ \ \underline{\hspace{0.5cm}} & C\ \ \underline{\hspace{0.5cm}} \\
\textbf{4.} Hvem har altid fri både lørdag og søndag? & A\ \ \underline{\hspace{0.5cm}} & B\ \ \underline{\hspace{0.5cm}} & C\ \ \underline{\hspace{0.5cm}} \\
\textbf{5.} Hvem bruger længst tid på at komme på arbejde? & A\ \ \underline{\hspace{0.5cm}} & B\ \ \underline{\hspace{0.5cm}} & C\ \ \underline{\hspace{0.5cm}} \\
\textbf{6.} Hvem vil gerne arbejde færre timer i fremtiden? & A\ \ \underline{\hspace{0.5cm}} & B\ \ \underline{\hspace{0.5cm}} & C\ \ \underline{\hspace{0.5cm}} \\
\textbf{7.} Hvem var uden job i en periode før det nuværende arbejde? & A\ \ \underline{\hspace{0.5cm}} & B\ \ \underline{\hspace{0.5cm}} & C\ \ \underline{\hspace{0.5cm}} \\
\bottomrule
\end{tabularx}
\end{center}

\vspace{0.2em}
\begin{tcolorbox}[enhanced, breakable, colback=NCDBgGray, colframe=NCDGray, arc=2pt, boxrule=0.7pt, borderline west={3pt}{0pt}{NCDNavy}, left=10pt, right=10pt, top=5pt, bottom=5pt]
\textbf{\color{NCDNavy}A. Freja}\\[2pt]
Jeg er sygeplejerske på et stort hospital i Odense, hvor jeg begyndte for seks år siden. Før jeg kom her, havde jeg arbejdet to år på et plejehjem, så jeg skiftede direkte fra det ene job til det andet. Jeg arbejder på fuldtid, altså 37 timer om ugen, med både dagvagter og aftenvagter, og jeg arbejder også hver tredje weekend. Arbejdet her er travlere end på plejehjemmet, men også mere spændende. Jeg kan selvfølgelig ikke arbejde hjemmefra, når jeg passer patienter, så jeg cykler 15 minutter hver vej. Mine kolleger er søde, og vi hjælper altid hinanden. Næste år håber jeg at gå ned på 30 timer, fordi jeg gerne vil have mere overskud til mine venner.
\end{tcolorbox}

\begin{tcolorbox}[enhanced, breakable, colback=NCDBgGray, colframe=NCDGray, arc=2pt, boxrule=0.7pt, borderline west={3pt}{0pt}{NCDTeal}, left=10pt, right=10pt, top=5pt, bottom=5pt]
\textbf{\color{NCDTeal}B. Karim}\\[2pt]
Jeg arbejder som softwareudvikler i et it-firma i København. Jeg startede i firmaet for to år siden, lige efter at jeg havde studeret datalogi i fem år, og jeg blev faktisk ansat allerede inden min sidste eksamen. Min arbejdsuge er på 37 timer, men jeg har fleksible arbejdstider og arbejder hjemmefra hver tirsdag og torsdag. Jeg møder aldrig om aftenen eller i weekenden, så jeg har altid tid til sport om lørdagen. Mine kolleger er dygtige, og vi spiser frokost sammen i kantinen kl. 12. De dage, hvor jeg tager ind på kontoret, tager turen med S-toget næsten 45 minutter hver vej. Det er lidt langt, så selvom jeg gerne vil blive ved med at arbejde fuldtid, leder jeg nu efter et job tættere på min lejlighed.
\end{tcolorbox}

\begin{tcolorbox}[enhanced, breakable, colback=NCDBgGray, colframe=NCDGray, arc=2pt, boxrule=0.7pt, borderline west={3pt}{0pt}{NCDRed}, left=10pt, right=10pt, top=5pt, bottom=5pt]
\textbf{\color{NCDRed}C. Lotte}\\[2pt]
Jeg er ekspedient i et lille bageri i Aarhus, hvor jeg først fik job for fire måneder siden. Før jeg begyndte her, havde jeg faktisk gået hjemme uden arbejde i otte måneder, fordi min gamle butik desværre lukkede. Jeg har valgt kun at arbejde 28 timer om ugen, så jeg har mere tid til min yngste søn og min familie om eftermiddagen. Jeg møder meget tidligt klokken halv seks fra tirsdag til lørdag og holder fri søndag og mandag. Vi har altid travlt om morgenen, så hjemmearbejde er naturligvis umuligt for mig. Heldigvis bor jeg kun fem minutters gang fra bageriet. Jeg tjener lidt mindre end før, men stemningen her er varm, og jeg vil ikke ændre noget foreløbig.
\end{tcolorbox}

% =====================================================================
% Del 2: Skrivning
% =====================================================================
\ncdsubsec{Del 2: Skrivning（书面表达 · minimum 70 ord）}

\subsection*{Opgave 1（官方必考私人邮件回复）}

\textbf{Situation:} Du har fået en mail fra din ven, Emil.\\
\textbf{Opgave:} Læs e-mailen, og skriv et svar til Emil. Du skal svare på alle Emils spørgsmål. Du skal skrive \textbf{minimum 70 ord}.

\begin{emailbox}{来信｜Fra: Emil}
Hej\\
Jeg har hørt, at du har fået et nyt arbejde! Hvad arbejder du med, og hvor ligger det? Hvor mange timer arbejder du om ugen? Hvordan er dine kolleger og chefen? Og hvad kan du bedst lide ved jobbet?\\[4pt]
Mange hilsner\\
Emil
\end{emailbox}

\begin{tcolorbox}[enhanced, colback=white, colframe=NCDGray, arc=2pt, boxrule=0.8pt, left=10pt, right=10pt, top=7pt, bottom=7pt]
\textbf{Til:} Emil \qquad\qquad \textbf{Emne:} Nyt arbejde\\[6pt]
\rule{\linewidth}{0.4pt}\\[9pt]
\rule{\linewidth}{0.4pt}\\[9pt]
\rule{\linewidth}{0.4pt}\\[9pt]
\rule{\linewidth}{0.4pt}\\[9pt]
\rule{\linewidth}{0.4pt}\\[9pt]
\rule{\linewidth}{0.4pt}\\[9pt]
\rule{\linewidth}{0.4pt}\\[9pt]
\rule{\linewidth}{0.4pt}
\end{tcolorbox}

\subsection*{Bonus Opgave 1B（半正式实用邮件训练：物业报修）}

\begin{tipbox}{任务说明与交际提示（Semiformel e-mail til viceværten）}
\small
\textbf{Situation:} Du bor i en lejlighed, hvor der er opstået to praktiske problemer: Radiatoren i din stue har ikke virket siden i onsdags, og lyset i det låste cykelskur virker heller ikke, så det er helt mørkt om morgenen.\\
\textbf{Opgave:} Skriv en høflig e-mail (\textbf{minimum 70 ord}) til bygningens vicevært, \textbf{Henrik}. Beskriv begge problemer tydeligt, og spørg høfligt, hvornår han kan komme og kigge på dem.\\[3pt]
\textbf{\color{NCDRed}[核心语法提示]} 在丹麦半正式沟通（如联系物业管理员 \textit{vicevært}、老师或邻居）中，\textbf{统一使用第二人称 \textit{du / dig / din / dit}}（不可混用过时的敬称 \textit{De/Dem}），用 \textit{Hej Henrik} 开头、\textit{Venlig hilsen} 结尾，并主动用礼貌过去时（\textit{Kunne du måske...? / Jeg ville sætte stor pris på, hvis...}）与时间介词 \textit{siden i onsdags}。
\end{tipbox}

\begin{tcolorbox}[enhanced, colback=white, colframe=NCDGray, arc=2pt, boxrule=0.8pt, left=10pt, right=10pt, top=7pt, bottom=7pt]
\textbf{Til:} Henrik (Vicevært) \qquad\qquad \textbf{Emne:} Radiator i stuen og lys i cykelskuret\\[6pt]
\rule{\linewidth}{0.4pt}\\[10pt]
\rule{\linewidth}{0.4pt}\\[10pt]
\rule{\linewidth}{0.4pt}\\[10pt]
\rule{\linewidth}{0.4pt}\\[10pt]
\rule{\linewidth}{0.4pt}\\[10pt]
\rule{\linewidth}{0.4pt}\\[10pt]
\rule{\linewidth}{0.4pt}\\[10pt]
\rule{\linewidth}{0.4pt}
\end{tcolorbox}

% =====================================================================
% Del 3: Mundtlig kommunikation
% =====================================================================
\ncdsubsec{Del 3: Mundtlig kommunikation（口语交际全真题卡）}

\subsection*{Mundtlig Opgave 1（Emneopgave med Mindmap · 陈述 1--2 分钟 + 追问 3--4 分钟）}

\begin{tipbox}{官方考规提醒（SIRI DU3 Modul 2 Opgave 1）}
\small
\ding{172} 考生考前准备两张主题 Mindmap（其中至少一张须属于 \textit{arbejde / praktik / studie / sprogskole} 领域），考场抽签选一张陈述 \textbf{1--2 分钟（Del 1）}，随后考官针对内容追问 \textbf{3--4 分钟（Del 2）}，\textbf{两部分均须通过}。\\
\ding{173} \textbf{铁律：Mindmap 纸上只能写关键词/短语（stikord），严禁写完整句子！}
\end{tipbox}

\begin{tcolorbox}[enhanced, breakable, colback=NCDBgBlue, colframe=NCDNavy, coltitle=white, fonttitle=\bfseries, title={Mindmap 2A：Min sprogskole（官方六大分支 · 仅限 Stikord）}, arc=3pt, boxrule=0.8pt, left=10pt, right=10pt, top=7pt, bottom=7pt]
\small
\renewcommand{\arraystretch}{1.35}
\begin{tabularx}{\linewidth}{@{} >{\raggedright\arraybackslash}p{8.0cm} >{\raggedright\arraybackslash}X @{}}
\textbf{\color{NCDNavy}1. sprogskole / by}\newline
\textit{Lærdansk · Aarhus · tæt på centrum · stor, moderne bygning}
&
\textbf{\color{NCDNavy}2. dage / tid}\newline
\textit{tirsdag \& torsdag formiddag · kl. 8.30--11.45 · to gange om ugen}
\\
\midrule
\textbf{\color{NCDNavy}3. kursister / lærer}\newline
\textit{14 kursister · mange lande · sød, tålmodig lærer (Mette)}
&
\textbf{\color{NCDNavy}4. aktiviteter}\newline
\textit{samtale i par · grammatik · lytteøvelser · små dialoger}
\\
\midrule
\textbf{\color{NCDNavy}5. pauser}\newline
\textit{kl. 10.00 (15 minutter) · kaffe i kantinen · tale dansk sammen}
&
\textbf{\color{NCDNavy}6. transport til sprogskolen}\newline
\textit{cykel (20 min.) · frisk luft · bussen i dårligt vejr}
\\
\end{tabularx}
\vspace{0.4em}
\hrule
\vspace{0.4em}
\textbf{\color{NCDRed}Del 2 考官追问模拟提示（Uddybende spørgsmål · 3--4 min.）：}
\begin{enumerate}
  \item \textit{Du siger, at I laver mange forskellige aktiviteter i timerne. Hvad kan du bedst lide at lave?}
  \item \textit{Hvad plejer I at snakke om i pausen i kantinen?}
  \item \textit{Har du fået nogle gode venner på dit hold? Hvor kommer de fra?}
  \item \textit{Hvad synes du er det sværeste ved at lære dansk indtil nu?}
\end{enumerate}
\end{tcolorbox}

\begin{tcolorbox}[enhanced, breakable, colback=NCDBgTeal, colframe=NCDTeal, coltitle=white, fonttitle=\bfseries, title={Mindmap 2B：Min fritidsaktivitet: Løbeklub / Fitness（官方六大分支 · 仅限 Stikord）}, arc=3pt, boxrule=0.8pt, left=10pt, right=10pt, top=7pt, bottom=7pt]
\small
\renewcommand{\arraystretch}{1.35}
\begin{tabularx}{\linewidth}{@{} >{\raggedright\arraybackslash}p{8.0cm} >{\raggedright\arraybackslash}X @{}}
\textbf{\color{NCDTeal}1. sted / by}\newline
\textit{Marselisborg Løbeklub · Aarhus · ved skoven og stranden}
&
\textbf{\color{NCDTeal}2. dage / tid}\newline
\textit{mandag \& onsdag aften · kl. 17.30--18.45 · nogle lørdage}
\\
\midrule
\textbf{\color{NCDTeal}3. pris / gratis}\newline
\textit{150 kr. om måneden · første måned gratis · billigere end fitnesscenter}
&
\textbf{\color{NCDTeal}4. alene / sammen med}\newline
\textit{sammen med 20 medlemmer · forskellige aldre · god ven (Jonas)}
\\
\midrule
\textbf{\color{NCDTeal}5. aktiviteter}\newline
\textit{opvarmning · løbe 5--8 km · strække ud · varm te bagefter}
&
\textbf{\color{NCDTeal}6. transport til aktiviteten}\newline
\textit{på cykel (10 min.) · letbanen i regnvejr}
\\
\end{tabularx}
\vspace{0.4em}
\hrule
\vspace{0.4em}
\textbf{\color{NCDRed}Del 2 考官追问模拟提示（Uddybende spørgsmål · 3--4 min.）：}
\begin{enumerate}
  \item \textit{Hvorfor valgte du at starte i en løbeklub i stedet for at løbe alene?}
  \item \textit{Du siger, at du også løber nogle lørdag formiddage. Hvad laver I efter træningen?}
  \item \textit{Hvordan har du det i kroppen efter en løbetur om aftenen?}
  \item \textit{Løber du også udenfor om vinteren, når det er mørkt og koldt?}
\end{enumerate}
\end{tcolorbox}

\subsection*{Mundtlig Opgave 2（Informationskløft · Case: \textit{Lukas} · 满分 6 分，4--6 分及格）}

\begin{tipbox}{考场规则提醒（DU3 Modul 2 Opgave 2: Informationskløft）}
\small
两名考生面对面坐好，互不偷看对方题卡。考生 A 抽到 \textbf{Kursist A Ark}，须用\textbf{第三人称（\textit{han / Lukas}）}针对 2 个主题下的全部 6 个分类格提出 6 个聚焦问题（可用直接问句或礼貌间接问句）。\textbf{切勿对右上角人物头像提问！} 考生 B 仅根据 \textbf{Kursist B Ark} 的图片线索作答；双方不得互相提示如何发问。\textbf{仅提问方在此题计分}：每个合格问题得 1 分，满分 6 分，\textbf{4--6 分及格}。随后 A/B 互换角色并抽取新案例。
\end{tipbox}

\begin{tcolorbox}[enhanced, breakable, colback=NCDBgGray, colframe=NCDNavy, coltitle=white, fonttitle=\bfseries, title={Kursist A Ark --- Case: Lukas（提问方题卡 · 含 6 个视觉问号格）}, arc=2pt, boxrule=0.8pt, left=10pt, right=10pt, top=7pt, bottom=7pt]
\small
\textbf{人物提示（Portrætfoto øverst til højre）：} \textit{Lukas}（年轻男大学生微笑半身照 --- 仅用于确认人物名字与性别，\textbf{不要对此照片提问}）。

\vspace{0.4em}
\renewcommand{\arraystretch}{1.35}
\begin{tabularx}{\linewidth}{@{} >{\raggedright\arraybackslash}p{8.0cm} >{\raggedright\arraybackslash}X @{}}
\toprule
\textbf{\color{NCDNavy}Emne 1: Lukas læser på universitetet} & \textbf{\color{NCDNavy}Emne 2: Lukas’ bolig} \\
\midrule
\textbf{1. studie / fag}\newline
[图片线索：大学阶梯教室与书本图标] \quad \textbf{\color{NCDRed}?}
&
\textbf{4. sted / by}\newline
[图片线索：城市街区与地图定位图标] \quad \textbf{\color{NCDRed}?}
\\
\midrule
\textbf{2. dage / tid}\newline
[图片线索：周历表格与时钟指针图标] \quad \textbf{\color{NCDRed}?}
&
\textbf{5. boligtype / værelser}\newline
[图片线索：住宅建筑与房间平面图标] \quad \textbf{\color{NCDRed}?}
\\
\midrule
\textbf{3. transport til universitetet}\newline
[图片线索：通勤路线与交通问号图标] \quad \textbf{\color{NCDRed}?}
&
\textbf{6. alene / sammen med}\newline
[图片线索：住房内人影轮廓对比图标] \quad \textbf{\color{NCDRed}?}
\\
\bottomrule
\end{tabularx}
\end{tcolorbox}

\begin{tcolorbox}[enhanced, breakable, colback=NCDBgGray, colframe=NCDTeal, coltitle=white, fonttitle=\bfseries, title={Kursist B Ark --- Case: Lukas（回答方线索卡 · 供 B 看图回答）}, arc=2pt, boxrule=0.8pt, left=10pt, right=10pt, top=7pt, bottom=7pt]
\small
\renewcommand{\arraystretch}{1.35}
\begin{tabularx}{\linewidth}{@{} >{\raggedright\arraybackslash}p{8.0cm} >{\raggedright\arraybackslash}X @{}}
\toprule
\textbf{\color{NCDTeal}Emne 1: Lukas læser på universitetet} & \textbf{\color{NCDTeal}Emne 2: Lukas’ bolig} \\
\midrule
\textbf{1.} [图片线索：显微镜与《Biologi》教材 · 校徽 Aarhus Universitet]\newline
$\rightarrow$ \textit{Han læser biologi på Aarhus Universitet.}
&
\textbf{4.} [图片线索：路牌“Aarhus V” · 森林边缘位置图标]\newline
$\rightarrow$ \textit{Han bor i Aarhus Vest tæt på skoven.}
\\
\midrule
\textbf{2.} [图片线索：周一至周五日程表 09:00--15:00 · 周三小组讨论]\newline
$\rightarrow$ \textit{Mandag til fredag kl. 9.00--15.00 (læsegruppe om onsdagen).}
&
\textbf{5.} [图片线索：学生宿舍楼 Kollegium · 18 m$^2$ 单间带独立卫浴]\newline
$\rightarrow$ \textit{På et kollegium (18 m$^2$ værelse med eget bad).}
\\
\midrule
\textbf{3.} [图片线索：轻轨列车 Letbanen · 计时器 20 min.]\newline
$\rightarrow$ \textit{Han tager letbanen, og det tager cirka 20 minutter.}
&
\textbf{6.} [图片线索：6 名大学生在共用厨房一起做饭]\newline
$\rightarrow$ \textit{Alene på sit værelse, men deler køkken med 5 andre studerende.}
\\
\bottomrule
\end{tabularx}
\end{tcolorbox}

% =====================================================================
% Del 4: Retteark, Detaljeret Analyse & Højscore-Modeller
% =====================================================================
\ncdsubsec{Del 4: Retteark, Detaljeret Analyse \& Højscore-Modeller（评分卡、解析与高分范文）}

% ---------------------------------------------------------------------
\subsection*{4.1 官方 Retteark 评分卡（Læsning 24 point + Mundtlig Opgave 2 6 point）}

\begin{grammatikbox}{Læsning（24 point · Bestået: 16--24 point · Ikke bestået: 0--15 point）}
\small
\renewcommand{\arraystretch}{1.24}
\begin{tabularx}{\linewidth}{@{} >{\bfseries\color{NCDNavy}\raggedright\arraybackslash}p{3.0cm} >{\raggedright\arraybackslash}X >{\centering\arraybackslash}p{2.0cm} @{}}
\toprule
\textbf{题号与分值} & \textbf{官方标准答案（Rettenøgle）} & \textbf{得分记录} \\
\midrule
\textbf{Opgave 1} (6 pt) &
\textbf{0.} A \textit{[Eksempel]} \quad|\quad \textbf{1.} \textbf{D} \quad·\quad \textbf{2.} \textbf{G} \quad·\quad \textbf{3.} \textbf{B} \quad·\quad \textbf{4.} \textbf{I} \quad·\quad \textbf{5.} \textbf{E} \quad·\quad \textbf{6.} \textbf{C}\newline
（多余干扰项：\textbf{F}、\textbf{H}） & \underline{\hspace{0.7cm}} / 6 \\
\midrule
\textbf{Opgave 2} (5 pt) &
每段须精准划去下列矛盾/离题句：\newline
\textbf{1.} \textit{Det var midt om natten, og stationen var helt tom.}\newline
\textbf{2.} \textit{Derfor åbnede han straks sin egen computer på bordet og rettede de sidste billeder.}\newline
\textbf{3.} \textit{Solen skinnede fra en helt blå himmel, og alle gik rundt i T-shirt.}\newline
\textbf{4.} \textit{Sofie gik straks sin vej og lukkede døren, fordi hun aldrig holdt jobsamtaler om tirsdagen.}\newline
\textbf{5.} \textit{Han var så vred over den dårlige samtale, at han aldrig ville besøge Odense igen.} & \underline{\hspace{0.7cm}} / 5 \\
\midrule
\textbf{Opgave 3} (6 pt) &
\textbf{0.} \textit{alle [Eksempel]} \quad|\quad \textbf{1.} \textbf{egne} \quad·\quad \textbf{2.} \textbf{end} \quad·\quad \textbf{3.} \textbf{hvor} \quad·\quad \textbf{4.} \textbf{bo} \quad·\quad \textbf{5.} \textbf{indtil} \quad·\quad \textbf{6.} \textbf{når}\newline
（4 个多余词：\textit{aldrig, havde, sælge, penge}） & \underline{\hspace{0.7cm}} / 6 \\
\midrule
\textbf{Opgave 4} (7 pt) &
\textbf{0.} A \textit{[Eksempel]} \quad|\quad \textbf{1.} \textbf{B} \quad·\quad \textbf{2.} \textbf{C} \quad·\quad \textbf{3.} \textbf{A} \quad·\quad \textbf{4.} \textbf{B} \quad·\quad \textbf{5.} \textbf{B} \quad·\quad \textbf{6.} \textbf{A} \quad·\quad \textbf{7.} \textbf{C} & \underline{\hspace{0.7cm}} / 7 \\
\midrule
\textbf{Læsning 总计} & \textbf{及格线：16--24 分}（建议冲刺目标：20+ 分） & \textbf{\underline{\hspace{0.7cm}} / 24} \\
\bottomrule
\end{tabularx}
\end{grammatikbox}

\begin{tipbox}{Skrivning \& Mundtlig 考官评分自检核对表}
\small
\begin{itemize}
  \item \textbf{Del 2 Skrivning}：字数是否达到 \textbf{70 词以上}？是否完整回答了来信中全部 \textbf{4 个问题}（或 Bonus 1B 的两项报修+礼貌询问时间）？是否检查了从句否定词位置（\textit{ikke} 在动词前）、V2 倒装以及过去完成时/时间介词准确性？
  \item \textbf{Del 3 Mundtlig Opgave 1}：Mindmap 是否只写 \textit{stikord}（未写整句）？陈述是否自然持续 1--2 分钟？Del 2 能否接住考官的追问并展开解释？
  \item \textbf{Del 3 Mundtlig Opgave 2（满分 6 分，4--6 分及格）}：6 个格子是否全部使用\textbf{第三人称（\textit{han / Lukas}）}且聚焦于格子提示词？每个合格聚焦问题计 1 分（\textbf{\underline{\hspace{0.6cm}} / 6}）。
\end{itemize}
\end{tipbox}

% ---------------------------------------------------------------------
\subsection*{4.2 阅读 Opgave 1--4 逐题同义转述解析、全文精译与核心词表}

\subsubsection*{一、Opgave 1 逐题同义转述与干扰项拆解 + 全文中文精译}

\begin{transbox}
\small
\textbf{\color{NCDTeal}[Personerne 0--6 中文译文与精准定位]}\\
\textbf{0. 梅特和亨里克（例题 $\rightarrow$ A）}：梅特和亨里克有两个小孩和一只大狗。他们想在西日德兰海滩附近租一周度假屋（\textit{sommerhus}），重要的是自己能做饭，且狗能在外面花园活动。\\
$\rightarrow$ \textbf{同义转述}：\textit{selv kan lave mad + tæt på stranden + hunden i en have} $\leftrightarrow$ \textbf{A} 中的 \textit{400 meter fra stranden, fuldt udstyret køkken, stor indhegnet have hvor børn og hunde kan lege frit}。

\textbf{1. 纳迪亚和五岁儿子 $\rightarrow$ \dabold{D}}：喜欢动物，想体验五天不一样的夏令营，儿子梦想近距离看山羊或母鸡、早上帮忙捡鸡蛋；有私家车，不必在城里。\\
$\rightarrow$ \textbf{同义转述}：\textit{se geder eller høns + samle æg + egen bil} $\leftrightarrow$ \textbf{D（Sønderby Bondegårdsferie）}中的 \textit{fodre geder, kaniner og høns + samle friske æg + I skal have egen bil}（住 3--7 天满足五天需求）。

\textbf{2. 莫恩斯和皮娅（58 岁）$\rightarrow$ \dabold{G}}：想在丹麦岛屿放松度假四天，懒得自己做饭（希望含早晚餐），喜欢在风景中安静骑车。\\
$\rightarrow$ \textbf{同义转述}：\textit{på en dansk ø + både morgenmad og aftensmad + cykle} $\leftrightarrow$ \textbf{G（Ærø Badehotel \& Cykelferie）}中的 \textit{på den smukke ø Ærø + morgenbuffet og toretters aftensmad + gratis lån af elcykler}。

\textbf{3. 卡斯珀与两位 22 岁同学 $\rightarrow$ \dabold{B}}：去北日德兰便宜周末游，想冲浪并认识年轻人；没车（乘火车和公交），需要交通方便且有公共厨房做饭。\\
$\rightarrow$ \textbf{同义转述}：\textit{billig + Nordjylland + surfe + møde andre unge + offentlig transport + fælleskøkken} $\leftrightarrow$ \textbf{B（Klitmøller Vandrerhjem）}中的 \textit{Billig overnatning i Thy i Nordjylland + surfsteder + ti minutters gang fra bussen + fælleskøkken + møder andre rejsende}。

\textbf{4. 记者丹尼尔 $\rightarrow$ \dabold{I}}：独自离家五天，在海边安静环境中写新书，需要好书桌和高速网络上午工作、下午海边散步。\\
$\rightarrow$ \textbf{同义转述}：\textit{alene + fem dage + ved havet + godt skrivebord + hurtigt internet} $\leftrightarrow$ \textbf{I（Bornholm Kysthytter）}中的 \textit{til én eller to gæster + ved Bornholms klipper og hav + godt arbejdsbord og lynhurtigt fibernet}（排除了毫无 Wi-Fi 的 \textbf{C}）。

\textbf{5. 比尔吉特（68 岁）与坐轮椅的丈夫 $\rightarrow$ \dabold{E}}：去大城市度过长周末，不用车就能逛博物馆、吃餐厅，住宿须靠近车站且有电梯。\\
$\rightarrow$ \textbf{同义转述}：\textit{kørestol + større by + tæt på stationen + elevator} $\leftrightarrow$ \textbf{E（Byhavnen Lejlighedshotel）}中的 \textit{midt i Aarhus + fem minutters gang fra hovedbanegården + bred elevator + indrettet til kørestolsbrugere}（直接排除了老建筑无电梯的 \textbf{G}）。

\textbf{6. 丽娜和托马斯 $\rightarrow$ \dabold{C}}：只想两人过一个彻底远离屏幕、工作和网络的安静周末；喜欢睡在大自然里但不要硬防潮垫，要真床和送到门口的早餐。\\
$\rightarrow$ \textbf{同义转述}：\textit{uden skærme og internet + sove ude i naturen + rigtig seng + morgenmad bragt til døren} $\leftrightarrow$ \textbf{C（Lilleø Naturshelters \& Glamping）}中的 \textit{hverken wi-fi eller tv + luksustelte i en stille skov + rigtige senge + leverer en kurv med frisk morgenmad til teltet}。

\vspace{0.35em}
\textbf{\color{NCDRed}[两大干扰项拆解]}\\
· \textbf{F（Solsikken 家庭露营地）}：位于法尔斯特岛，虽有户外泳池和儿童迪斯科，但必须自带帐篷/房车、全区覆盖 Wi-Fi 且\textbf{没有公共厨房}，既不符合 3（无厨房、无冲浪），也不符合 6（要真床、无网且安静）。\\
· \textbf{H（锡尔克堡湖区木屋）}：可供 6--10 人团体入住且\textbf{最短整租 7 晚}、无餐饮、阁楼楼梯陡峭，与 1--6 题的人数（1--3 人）、天数（2--5 天）或无障碍要求均冲突。
\end{transbox}

\begin{transbox}
\small
\textbf{\color{NCDTeal}[Annoncer A--I 广告全文精译]}\\
\textbf{A.} 在西日德兰布洛旺德距海滩仅 400 米处租一栋温馨度假屋！三间卧室、设备齐全的厨房及可供孩子和狗自由玩耍的大片围栏花园。周租 5,400 克朗。\\
\textbf{B.} 北日德兰蒂地区平价住宿，步行十分钟即到公交站与丹麦最佳冲浪点之一！为年轻活跃旅客提供每晚 240 克朗起的简朴客房，设大公共厨房与温馨客厅，禁带宠物。\\
\textbf{C.} 渴望远离手机与邮件的宁静吗？我们的豪华帐篷坐落于静谧森林中，仅限成年伴侣入住（禁带儿童及宠物）。睡真床配暖被，每日早晨送新鲜早餐篮至帐篷。无 Wi-Fi 及电视。\\
\textbf{D.} 在菲英岛有机农场上的温馨度假公寓体验乡村生活！每天早上 8:30 孩子们可帮忙喂山羊、兔子和鸡，并捡新鲜鸡蛋作早餐。带小厨房和大游乐场；无公交直达，须自驾（可订 3--7 天）。\\
\textbf{E.} 位于奥胡斯市中心的舒适公寓式酒店，步行五分钟即达中央火车站！全楼配宽敞电梯，多套房型专设轮椅无障碍卫浴。门外即是博物馆、咖啡馆与餐厅。禁带宠物。\\
\textbf{F.} 带着全家去法尔斯特岛玛丽艾吕斯特海滩享受热闹露营吧！须自带帐篷或房车，七月每晚有露天泳池和儿童迪斯科；全区高速 Wi-Fi，带小卖部但无公共厨房，欢迎带狗。\\
\textbf{G.} 在美丽的艾尔岛经典海滨酒店享受悠闲假期！房价含每日自助早餐及餐厅两道式晚餐，并免费借用电动自行车。适合成年人放松；客房位于无电梯老建筑二层（丹麦 1. sal）。可订 2--7 晚。\\
\textbf{H.} 适合 6--10 人大家庭或朋友团体的宽敞原木小屋！靠近锡尔克堡森林，仅按整周出租（周六至周六至少 7 晚）。四间卧室、柴火炉及皮划艇通道，不供餐，阁楼楼梯陡峭，建议驾车。\\
\textbf{I.} 博恩霍尔姆岛悬崖海边专供 1--2 人入住的安静小木屋！每栋配小厨房、舒适办公桌及超快光纤，可专心工作或海边漫步。勒讷港有便利公交，全年营业（可住 3--14 晚）。无餐厅、禁养宠物。
\end{transbox}

\subsubsection*{二、Opgave 2 逻辑排雷解析与全文中文精译}

\begin{transbox}
\small
\textbf{\color{NCDTeal}[全文中文精译与矛盾句解析]}\\
\textbf{(0)} 周二早晨安德斯起得很早，因为他要去欧登塞一家大公司参加求职面试。前一晚他就已经熨好了最体面的衬衫，并把简历装进黑色文件夹。\textcolor{NCDRed}{[例题划掉：安德斯直接又回床睡到下午三点。]} 接着他匆匆吃完早餐，提前骑车前往奥胡斯中央火车站。\\
\textbf{(1)} 走上站台时，已有许多忙碌的通勤者在晨寒中候车。安德斯三天前就买好了车票，于是从容地在静音车厢找了个座位。去欧登塞的火车准点在 07:15 发车。\dabold{[划掉：当时正值半夜，车站里空无一人。]}（\textit{解析}：与上文“周二早晨 07:15”和“许多忙碌的通勤者”直接时间矛盾。）他靠窗坐下，最后复习了一遍面试笔记。\\
\textbf{(2)} 车程刚走一半，安德斯突然发现自己把笔记本电脑忘在家里厨房餐桌上了。偏偏他昨晚为了做一份重要的面试演示文稿熬到很晚。安德斯顿时非常紧张，因为所有文件只存在那台落下的电脑里。\dabold{[划掉：于是他立刻打开桌上的自己的电脑，修改了最后几张图片。]}（\textit{解析}：前句刚交代电脑忘在奥胡斯家中、文件只在忘带的电脑里，不可能当场打开修改。）幸好他迅速联系上奥胡斯的女友/女性朋友，对方把演示文稿发到了他的邮箱。\\
\textbf{(3)} 火车抵达欧登塞火车站时，突然下起了大雨。安德斯没带伞，于是宁愿坐公交车也不愿冒雨步行二十分钟。在公交站，一位好心的女士帮他找到了去往商务园最快的线路。\dabold{[划掉：太阳从湛蓝的天空中照耀着，所有人都穿着 T 恤走来走去。]}（\textit{解析}：与本段“突然下起大雨、没带伞、不想冒雨走路”气候矛盾。）安德斯上了公交车，在面试前十分钟赶到了公司。\\
\textbf{(4)} 在前台，主管索菲热情地接待了安德斯，并给他端来一杯咖啡。她早已读过他的求职信，觉得他过去的经历非常契合岗位。五分钟后，他们走进一间明亮的会议室，另一位同事也已在座。\dabold{[划掉：索菲立刻转身离开并关上了门，因为她周二从不主持求职面试。]}（\textit{解析}：与索菲热情接待、进入会议室以及下文顺利开展面试彻底矛盾。）安德斯很快在屏幕上打开了演示文稿，面试比他担心的还要顺利。\\
\textbf{(5)} 会议结束时，索菲感谢了这场精彩的面试，并承诺周五中午前给予答复。返回车站的路上，安德斯不再紧张，虽然疲惫却大大松了一口气。他立刻给在奥胡斯帮过自己的朋友发了一条开心的消息，再次感谢她白天的相助。\dabold{[划掉：他对这场糟糕的面试气愤不已，发誓再也不想来欧登塞了。]}（\textit{解析}：与上文“比担心的还要顺利”“大大松了一口气”“开心的消息”完全相反。）坐上回奥胡斯的火车后，他买了一个牛角包，惬意地享受安静的归途。
\end{transbox}

\subsubsection*{三、Opgave 3 完形选词语法解析与全文精译}

\begin{transbox}
\small
\textbf{\color{NCDTeal}[全文中文精译]}\\
在许多丹麦城市的边缘地带，分布着一片片绿色花园，里面坐落着五颜六色的小木屋。这就是存在已逾百年的丹麦著名“市民农园（\textit{kolonihaver}）”。起初，这些花园是专为城市工人设计的，好让他们能种植属于自己的蔬菜并呼吸新鲜空气。如今，市民农园比以往任何时候都更受欢迎，尤其是在没有阳台的小公寓里居住的年轻家庭当中。夏天里，园主们几乎每个周末都在园子里度过，在那里烧烤、读书并与邻居闲聊。许多农园协会有着严格规定：由于木屋并非为寒冷的丹麦冬天而建，人们只能在十月底之前住在木屋里。一年中其余时间木屋都空着，但当春天再次归来时，花园又会重新充满鲜花、孩子和咖啡香气。

\vspace{0.35em}
\textbf{\color{NCDNavy}[逐空词法与干扰项排除解析]}\\
· \textbf{(1) egne}（复数形容词“自己的”）：位于所有格代词 \textit{deres} 与复数名词 \textit{grøntsager} 之间，只有复数形 \textit{egne} 符合词性与语义。\\
· \textbf{(2) end}（比较连词“比”）：紧接比较级 \textit{mere populære}，构成固定搭配 \textit{mere populære end nogensinde før}（比以往任何时候都更受欢迎）。\\
· \textbf{(3) hvor}（关系副词“在那里”）：引导定语从句修饰地点名词 \textit{haven}，后面接完整主谓结构 \textit{de griller}（不可填 \textit{når}，因后面无主句承接）。\\
· \textbf{(4) bo}（不定式动词“居住”）：情态动词 \textit{må kun} 后接动词原形，且后面接地点状语 \textit{i sit hus}（干扰项 \textit{sælge}“售卖”为及物动词且与“不能过冬”逻辑完全无关）。\\
· \textbf{(5) indtil}（时间介词“直到……为止”）：接时间终点 \textit{slutningen af oktober}，精准呼应后句“因为房子不是为冬天建的”。\\
· \textbf{(6) når}（时间从句连词“每当/当……时”）：用于一般现在时描述每年春天循环发生的现象，连接 \textit{men} 引出的转折分句。\\
· \textbf{4 个干扰词排除}：\textit{aldrig}（副词，任何空均不合语法）、\textit{havde}（过去时助动词，文本全篇为客观一般现在时事实陈述且空(4)须用原形）、\textit{sælge}（须接宾语且语义荒谬）、\textit{penge}（名词，全篇无任何名词空位）。
\end{transbox}

\subsubsection*{四、Opgave 4 人物对比同义转述解析与全文精译}

\begin{transbox}
\small
\textbf{\color{NCDTeal}[全文中文精译与逐题定位]}\\
\textbf{A. 弗蕾娅}：我是欧登塞一家大型医院的护士，六年前来到这里工作。来这儿之前，我曾在一家养老院工作了两年，所以我是直接从一份工作换到另一份工作的。我干的是全职，也就是每周 37 小时，既有白班也有晚班，每三个周末还得值一次班。这里的工作比养老院更忙，但也更刺激有趣。照顾病人当然不可能居家办公，所以我单程骑车 15 分钟上下班。同事们都很好，大家总是互相帮忙。明年我希望把工时降到 30 小时，因为我想有更多精力陪陪朋友。\\
\textbf{B. 卡里姆}：我在哥本哈根一家 IT 公司担任软件开发员。两年前我入职这家公司，在此之前我刚学了五年计算机科学，实际上在最后一场考试前就已经拿到了录用通知。我每周工作 37 小时，但工作时间很灵活，每周二和周四都在家办公。我从不在晚上或周末上班，所以周六总有时间运动。我的同事都很专业，大家十二点一起在食堂吃午饭。需要去办公室的那几天，坐城铁 S-tog 单程将近 45 分钟。这有点远，所以虽然我想继续做全职，但我现在正在找离公寓更近的工作。\\
\textbf{C. 洛特}：我是奥胡斯一家小面包店的店员，四个月前才刚刚拿到这份工作。来这里之前，由于原来的老店不幸倒闭，我其实赋闲在家整整八个月。我选择每周只工作 28 小时，这样下午就有更多时间陪伴小儿子和家人。我周二到周六早上五点半就得早早到岗，周日和周一休息。清晨店里总是特别忙，居家办公对我而言自然是不可能的。幸好我住的地方离面包店步行只要五分钟。我的薪水比以前略少一点，但店里气氛很温暖，暂时我并不打算做任何改变。

\vspace{0.35em}
\textbf{\color{NCDNavy}[Opgave 4 同义转述考点矩阵]}\\
· \textbf{1 $\rightarrow$ B}：\textit{arbejde hjemme nogle dage om ugen} $\leftrightarrow$ \textit{arbejder hjemmefra hver tirsdag og torsdag}（A、C 均明言无法居家办公）。\\
· \textbf{2 $\rightarrow$ C}：\textit{arbejder på deltid} $\leftrightarrow$ \textit{kun at arbejde 28 timer om ugen}（A、B 均为 37 小时全职）。\\
· \textbf{3 $\rightarrow$ A}：\textit{det samme sted i længst tid} $\leftrightarrow$ \textbf{6 år}（A）vs 2 år（B）vs 4 måneder（C）（最高级对比）。\\
· \textbf{4 $\rightarrow$ B}：\textit{altid fri både lørdag og søndag} $\leftrightarrow$ \textit{møder aldrig... i weekenden}（A 每三周值周末，C 周六上班）。\\
· \textbf{5 $\rightarrow$ B}：\textit{bruger længst tid på at komme på arbejde} $\leftrightarrow$ \textbf{45 min.}（B）vs 15 min.（A）vs 5 min.（C）。\\
· \textbf{6 $\rightarrow$ A}：\textit{arbejde færre timer i fremtiden} $\leftrightarrow$ \textit{Næste år håber jeg at gå ned på 30 timer}（B 要继续全职，C 不想改变）。\\
· \textbf{7 $\rightarrow$ C}：\textit{uden job i en periode før det nuværende arbejde} $\leftrightarrow$ \textit{havde jeg faktisk gået hjemme uden arbejde i otte måneder}（过去完成时回溯；A 无缝跳槽，B 毕业前即入职）。
\end{transbox}

\subsubsection*{五、本卷核心高频词表（Gloser · 15 个核心词汇与完整屈折）}

\begin{glosertable}
名词 & en overnatning & \textit{(overnatningen, overnatninger, overnatningerne)} & 过夜，住宿 \\
名词 & en jobsamtale & \textit{(jobsamtalen, jobsamtaler, jobsamtalerne)} & 求职面试 \\
名词 & en pendler & \textit{(pendleren, pendlere, pendlerne)} & 通勤者 \\
名词 & en kolonihave & \textit{(kolonihaven, kolonihaver, kolonihaverne)} & 市民农园，小木屋菜园 \\
名词 & en forening & \textit{(foreningen, foreninger, foreningerne)} & 协会，社团，俱乐部 \\
名词 & en vicevært & \textit{(viceværten, viceværter, viceværterne)} & 物业管理员，楼管 \\
名词 & en radiator & \textit{(radiatoren, radiatorer, radiatorerne)} & 暖气片 \\
名词 & et cykelskur & \textit{(cykelskuret, cykelskure, cykelskurene)} & 自行车棚 \\
名词 & et kollegium & \textit{(kollegiet, kollegier, kollegierne)} & 学生宿舍（不规则变格） \\
\midrule
动词 [har] & at stryge & \textit{(stryger, strøg, har strøget)} & 熨烫（衣物）；划掉 \\
动词 [har] & at tilbringe & \textit{(tilbringer, tilbragte, har tilbragt)} & 度过（时间、周末） \\
动词 [har] & at dyrke & \textit{(dyrker, dyrkede, har dyrket)} & 种植；从事（运动、爱好） \\
动词 [har] & at reparere & \textit{(reparerer, reparerede, har repareret)} & 修理，维修 \\
\midrule
形容词 & travl & \textit{travl (travlt, travle; travlere, travlest)} & 忙碌的（注意比较级 \textit{travlere}） \\
形容词 & fleksibel & \textit{fleksibel (-t, fleksible; mere/mest fleksibel)} & 灵活的，弹性的 \\
\end{glosertable}

% ---------------------------------------------------------------------
\subsection*{4.3 Del 2 写作高分双档范文（A2 稳过版 + A2+ 冲刺版 + 物业报修）}

\begin{emailbox}{Opgave 1 范文一：A2 稳过版（共 87 词 · 简洁零失误）}
Hej Emil

Tak for din mail! Ja, jeg har fået et nyt deltidsjob på en lille café i Aarhus. Caféen ligger i centrum tæt på hovedbanegården. Jeg hjælper med at lave kaffe, smøre sandwich og rydde op ved bordene.

Jeg arbejder 15 timer om ugen, mest om eftermiddagen og hver anden lørdag. Mine kolleger og min chef er rigtig søde, og vi hjælper altid hinanden, når der er travlt. Det, jeg bedst kan lide ved jobbet, er at snakke dansk med kunderne hver dag.

Mange hilsner\\
Sofie
\end{emailbox}

\begin{emailbox}{Opgave 1 范文二：A2+ 高分示范版（共 112 词 · 亮出本卷语法靶点）}
Hej Emil

Tusind tak for din mail! Det er rigtigt, at jeg \da{for tre uger siden} fik et nyt studiejob i en hyggelig boghandel midt i Aarhus. \dabold{Selvom jeg ikke havde arbejdet} i en butik før, lærte jeg hurtigt arbejdet at kende.

Jeg arbejder 16 timer om ugen, normalt tirsdag og torsdag eftermiddag \da{indtil} kl. 18 samt hver anden lørdag. Mine opgaver er at betjene kunder og sætte nye bøger på plads. Min chef og de tre kolleger er meget hjælpsomme, selvom der nogle gange er \dabold{travlere end} på mit gamle job.

\dabold{Det bedste} ved jobbet er uden tvivl, at jeg taler dansk hele tiden og møder søde mennesker.

Mange hilsner\\
Sofie
\end{emailbox}

\begin{tipbox}{双档范文精评（A2 vs A2+ 拉分关键）}
\small
\ding{172} \textbf{全覆盖四问}：两篇均完整回答了“做什么工作及地点”、“每周工时”、“同事与老板评价”、“最喜欢的部分”全部四个问题。\\
\ding{173} \textbf{A2+ 语法升级点}：高分版精准嵌入了 \textbf{时间介词}（\textit{for tre uger siden} 配一般过去时 \textit{fik}；\textit{indtil kl. 18} 表时间终点）、\textbf{过去完成时 + 从句否定语序}（\textit{Selvom jeg \textbf{ikke havde arbejdet}...}，从句前移触发主句倒装 \textit{lærte jeg}），以及\textbf{比较级与最高级}（\textit{travlere end}、\textit{Det bedste}）。
\end{tipbox}

\begin{tcolorbox}[enhanced, colback=NCDBgGray, colframe=NCDGray, coltitle=white, fonttitle=\bfseries, title={Bonus Opgave 1B 示范：致物业 Henrik 的礼貌半正式报修信（共 102 词）}, arc=2pt, boxrule=0.8pt, borderline west={3pt}{0pt}{NCDRed}, left=10pt, right=10pt, top=7pt, bottom=7pt]
\textbf{Til:} Henrik (Vicevært) \qquad \textbf{Emne:} Radiator i stuen og lys i cykelskuret\\[4pt]
Hej Henrik

Jeg bor i lejlighed 2. th og skriver til \da{dig}, fordi vi desværre har to problemer i bygningen.

For det første har radiatoren i min stue \dabold{ikke virket siden i onsdags}. Den er helt kold, selvom jeg har skruet helt op for den, og lejligheden er blevet \dabold{koldere}. For det andet virker lyset i det låste cykelskur heller ikke, så det er meget mørkt, når jeg henter min cykel om morgenen.

\dabold{Kunne du måske} kigge forbi og reparere tingene i løbet af ugen? Jeg \dabold{ville sætte stor pris på}, hvis \da{du} kan fortælle mig, \dabold{hvornår du kommer}.

Venlig hilsen\\
Sofie
\end{tcolorbox}

% ---------------------------------------------------------------------
\subsection*{4.4 Del 3 口语全真示范逐字稿、核心 IPA 标注与信息差问句矩阵}

\begin{tcolorbox}[enhanced, breakable, colback=NCDBgBlue, colframe=NCDNavy, coltitle=white, fonttitle=\bfseries, title={Mindmap 2A 口语示范稿：Min sprogskole（130 词 · 约 1.5 分钟）}, arc=3pt, boxrule=0.8pt, left=10pt, right=10pt, top=7pt, bottom=7pt]
\small
Jeg går på \textbf{sprogskolen} \ipa{ˈsbʁɔʊ̯ˌsɡoːlən} Lærdansk i Aarhus, som ligger tæt på centrum i en stor, moderne bygning. Jeg har \textbf{undervisning} \ipa{ˈɔnɐˌviːˀsneŋ} to gange om ugen, nemlig tirsdag og torsdag formiddag fra klokken 8.30 til 11.45. På mit hold er vi fjorten \textbf{kursister} \ipa{kʰuɐ̯ˈsisdɐ} fra mange \textbf{forskellige} \ipa{fʌˈsɡɛlˀiə} lande, og vi har en rigtig sød og \textbf{tålmodig} \ipa{tˢʌlˈmoːˀði} lærer, der hedder Mette. I timerne laver vi mange forskellige \textbf{aktiviteter} \ipa{aɡtˢiviˈtˢeːˀdɐ}. Vi snakker meget sammen i par, øver grammatik og lytter til små \textbf{dialoger} \ipa{diaˈloːˀɐ}. Klokken ti har vi en \textbf{pause} \ipa{ˈpʰɑʊ̯sə} på femten minutter, hvor jeg næsten altid drikker kaffe i \textbf{kantinen} \ipa{kʰanˈtˢiːnən} og taler dansk med mine \textbf{klassekammerater} \ipa{ˈkʰlasəkʰamaˌʁɑːˀdɐ}. Jeg cykler til sprogskolen på tyve minutter, hvilket giver mig god motion og frisk luft før timerne. Når \textbf{vejret} \ipa{ˈvɛːˀɐð} bliver dårligere om vinteren, tager jeg dog bussen \textbf{i stedet for} \ipa{i ˈsdɛːˀðəð fʌ}.

\vspace{0.35em}
\textbf{\color{NCDTeal}[中文参考译文]}\\
我在奥胡斯的 Lærdansk 语言学校上课，学校靠近市中心，在一栋宽敞现代的大楼里。我每周上两次课，分别是周二和周四上午 8:30 到 11:45。我们班有十四名来自不同国家的学员，还有一位非常友善且有耐心的老师，名叫梅特。课堂上我们做很多不同的练习：结对对话、操练语法以及听小对话。十点整有十五分钟休息，我几乎总会在食堂喝杯咖啡，并用丹麦语和同学们聊天。我骑车去学校大概要二十分钟，课前正好能做些运动、呼吸新鲜空气。不过冬天天气变差时，我会改坐公交车。
\end{tcolorbox}

\begin{transbox}
\small
\textbf{\color{NCDNavy}[Mindmap 2A --- Del 2 考官追问高分应答示范]}\\
\textbf{Q1:} \textit{Du siger, at I laver mange forskellige aktiviteter i timerne. Hvad kan du bedst lide at lave?}\\
\textbf{A1:} \textit{Jeg kan bedst lide at øve dialoger i par, fordi vi får talt meget dansk og lærer nye ord, som vi også kan bruge ude i byen.}\\
\textbf{Q2:} \textit{Hvad plejer I at snakke om i pausen i kantinen?}\\
\textbf{A2:} \textit{Vi snakker tit om vores arbejde, vores familier og om vejret, og nogle gange spørger vi hinanden om lektierne til næste gang.}\\
\textbf{Q3:} \textit{Har du fået nogle gode venner på dit hold? Hvor kommer de fra?}\\
\textbf{A3:} \textit{Ja, jeg snakker rigtig godt med en kursist fra Brasilien og én fra Tyskland. Nogle gange drikker vi en kop kaffe sammen efter undervisningen.}\\
\textbf{Q4:} \textit{Hvad synes du er det sværeste ved at lære dansk indtil nu?}\\
\textbf{A4:} \textit{Udtalen er helt klart det sværeste for mig, fordi mange bogstaver ikke bliver udtalt, men det bliver nemmere, jo mere jeg lytter til dansk.}
\end{transbox}

\begin{tcolorbox}[enhanced, breakable, colback=NCDBgTeal, colframe=NCDTeal, coltitle=white, fonttitle=\bfseries, title={Mindmap 2B 口语示范稿：Min fritidsaktivitet: Løbeklub（130 词 · 约 1.5 分钟）}, arc=3pt, boxrule=0.8pt, left=10pt, right=10pt, top=7pt, bottom=7pt]
\small
I min \textbf{fritid} \ipa{ˈfʁiˌtˢiðˀ} går jeg i Marselisborg \textbf{Løbeklub} \ipa{ˈløːbəˌkʰlub} her i Aarhus, hvor vi mødes ved skoven og stranden. Vi træner hver mandag og onsdag aften fra klokken 17.30 til 18.45 samt nogle lørdag formiddage. Det koster kun 150 kroner om \textbf{måneden} \ipa{ˈmɔːnəðən}, og den første måned var helt gratis, så det er meget \textbf{billigere} \ipa{ˈbiliɐ} end et almindeligt fitnesscenter. Jeg løber aldrig alene, fordi vi er omkring tyve \textbf{medlemmer} \ipa{ˈmɛðˌlɛmˀɐ} i forskellige aldre på holdet, og min gode ven Jonas går der også. Når vi mødes, begynder vi med \textbf{opvarmning} \ipa{ˈʌbˌvɑːˀmneŋ}, og derefter løber vi mellem fem og otte \textbf{kilometer} \ipa{kʰiloˈmeːˀdɐ} i vores eget tempo. Bagefter strækker vi ud og drikker en kop varm te sammen. Det tager mig kun ti minutter at cykle hen til klubben, eller jeg tager \textbf{letbanen} \ipa{ˈlɛdˌbæːnən}, hvis det \textbf{regner} \ipa{ˈʁɑi̯nɐ} meget.

\vspace{0.35em}
\textbf{\color{NCDTeal}[中文参考译文]}\\
业余时间我在奥胡斯参加马塞利斯堡跑步俱乐部，我们在森林和海滩边集合。每周一和周三晚上 17:30 到 18:45 训练，有些周六上午也会跑。月费只要 150 克朗，第一个月还完全免费，比普通健身房便宜多了。我从不一个人跑，因为班里有大约二十位不同年龄的会员，而且我的好朋友约纳斯也参加。集合后我们先热身，接着按各自配速跑五到八公里。跑完大家拉伸肌肉，还一起喝杯热茶。我骑车去俱乐部只要十分钟，要是雨下得大我就坐轻轨。
\end{tcolorbox}

\begin{transbox}
\small
\textbf{\color{NCDTeal}[Mindmap 2B --- Del 2 考官追问高分应答示范]}\\
\textbf{Q1:} \textit{Hvorfor valgte du at starte i en løbeklub i stedet for at løbe alene?}\\
\textbf{A1:} \textit{Fordi det er meget sjovere at løbe sammen med andre. Når man har en fast aftale, kommer man også afsted, selvom man er træt efter arbejde.}\\
\textbf{Q2:} \textit{Du siger, at du også løber nogle lørdag formiddage. Hvad laver I efter træningen?}\\
\textbf{A2:} \textit{Om lørdagen har vi bedre tid, så vi sætter os tit på bænkene ved skoven og drikker varm te og spiser et stykke frugt sammen.}\\
\textbf{Q3:} \textit{Hvordan har du det i kroppen efter en løbetur om aftenen?}\\
\textbf{A3:} \textit{Jeg bliver lidt træt i benene, men jeg får meget mere energi i hovedet, og jeg sover meget bedre om natten.}\\
\textbf{Q4:} \textit{Løber du også udenfor om vinteren, når det er mørkt og koldt?}\\
\textbf{A4:} \textit{Ja, men så tager jeg varmt tøj, hue og en refleksvest på, og vi løber på veje med godt lys.}
\end{transbox}

\clearpage
\subsubsection*{Mundtlig Opgave 2（Informationskløft：\textit{Lukas}）双轨提问示范表与第三人称转述}

\begin{grammatikbox}{Lukas 6 格第三人称提问全表（直接满分句型 vs 间接不倒装高分句型 + B 卡回答）}
\small
\renewcommand{\arraystretch}{1.26}
\begin{tabularx}{\linewidth}{@{} >{\bfseries\raggedright\arraybackslash}p{2.1cm} >{\raggedright\arraybackslash}p{4.3cm} >{\raggedright\arraybackslash}p{4.8cm} >{\raggedright\arraybackslash}X @{}}
\toprule
\textbf{提示分类} & \textbf{直接问法（基础满分）} & \textbf{礼貌间接问法（高分不倒装）} & \textbf{Kursist B 回答（第三人称）} \\
\midrule
1. studie / fag\newline （专业与学校） &
\textit{Hvad læser Lukas, og hvor studerer han?} &
\textit{Kan du fortælle mig, \da{hvad} \dabold{Lukas læser}, og hvor han studerer?} &
\textit{Han læser biologi på Aarhus Universitet.} \\
\midrule
2. dage / tid\newline （上课时间） &
\textit{Hvornår har han undervisning på universitetet?} &
\textit{Jeg vil gerne vide, \da{hvornår} \dabold{han har} undervisning.} &
\textit{Fra mandag til fredag kl. 9.00--15.00, og om onsdagen har han læsegruppe.} \\
\midrule
3. transport\newline （上学交通） &
\textit{Hvordan kommer Lukas hen til universitetet?} &
\textit{Ved du, \da{hvordan} \dabold{han kommer} til universitetet, og hvor lang tid det tager?} &
\textit{Han tager letbanen, og det tager cirka 20 minutter.} \\
\midrule
4. sted / by\newline （居住地点） &
\textit{Hvor bor Lukas henne?} &
\textit{Kan du sige mig, \da{hvor} \dabold{Lukas bor}?} &
\textit{Han bor i Aarhus Vest tæt på skoven.} \\
\midrule
5. boligtype /\newline værelser &
\textit{Hvilken boligtype bor han i, og hvor stort er det?} &
\textit{Jeg vil gerne høre, \da{hvordan} \dabold{han bor}, og hvor stort hans værelse er.} &
\textit{Han bor på et kollegium på et 18 m$^2$ værelse med eget bad.} \\
\midrule
6. alene /\newline sammen med &
\textit{Bor han alene, eller bor han sammen med andre?} &
\textit{Ved du, \da{om} \dabold{Lukas bor} alene, eller om han deler med andre?} &
\textit{Han bor alene på sit værelse, men han deler køkken med fem andre studerende.} \\
\bottomrule
\end{tabularx}
\end{grammatikbox}

\begin{emailbox}{Mundtlig Opgave 2 口语汇总转述示范（第三人称 · 共 80 词）}
Lukas læser biologi på Aarhus Universitet, hvor han begyndte efter sommerferien. Han har normalt undervisning fra mandag til fredag mellem klokken ni og femten, men om onsdagen mødes han ofte med sin læsegruppe. Når han skal til universitetet, tager han letbanen, og turen tager cirka tyve minutter. Lukas bor i den vestlige del af Aarhus på et kollegium tæt på skoven. Han har sit eget værelse på atten kvadratmeter med eget badeværelse, men han deler køkken med fem andre studerende.
\end{emailbox}
